\section{What restart prices optimize}\label{sec:dual48}
For bounded nonnegative Borel prices $\lambda_t$ with $\lambda_T=0$, define $u_T=0$ and
\begin{equation}\label{eq:price48}
 u_t(x)=\min_a\left\{k_t(x,a)+\lambda_t(x)(d_t(x,a)-\varepsilon)
 +\beta P_t^a\!\left[u_{t+1}+\varepsilon\min\{\lambda_t(x),\lambda_{t+1}\}\right](x)\right\}.
\end{equation}
Every feasible common policy satisfies $C_t^p+\lambda_t(D_t^p-\varepsilon)\ge u_t$. Indeed, backward substitution leaves the nonnegative hinge
\[
 (\lambda_t(x)-\lambda_{t+1}(y))_+D_{t+1}^p(y)
 +(\lambda_{t+1}(y)-\lambda_t(x))_+(\varepsilon-D_{t+1}^p(y)).
\]
Together with $C_t^p\ge0$, this proves the lower bound
\begin{equation}\label{eq:clipped48}
 L(\lambda)=\int\max\{0,u_0\}\,d\nu\le\KR(\nu,\varepsilon).
\end{equation}
The same induction permits lower disadvantages in place of $d$, an observation used by the implemented witness-only method.

\subsection{An exact finite occupation interpretation}
For this subsection let $X=\{1,\ldots,n\}$. Write $P_{tia,j}$ for the transition probability and let $\mu$ be any nonnegative initial subprobability. For nonnegative variables $y_{tia}$ define $z_{ti}=\sum_a y_{tia}$ and $f_{tij}=\beta\sum_a y_{tia}P_{tia,j}$. Edges only run from date $t$ to date $t+1$, for $t<T-1$. Set missing edge sums to zero. Consider the linear program
\begin{align}
 \mathcal O(\mu)=\min_{y,\alpha,\gamma\ge0}\quad&\sum_{t,i,a}k_{tia}y_{tia},\label{eq:occupation48}\\
 z_{0i}=\mu_i,\quad z_{ti}&=\sum_h f_{t-1,hi}\quad(t\ge1),\nonumber\\
 \sum_a y_{tia}(\varepsilon-d_{tia})&\ge\sum_j\alpha_{tij}+\sum_h\gamma_{t-1,hi},\label{eq:budget48}\\
 \alpha_{tij}+\gamma_{tij}&\ge\varepsilon f_{tij}.\label{eq:edge48}
\end{align}
The variables $y$ are discounted occupations. The two nonnegative allocations on an incoming edge need not encode the same successor regret as allocations from other incoming states.

\begin{theorem}[Price supremum and common-budget completion]\label{thm:flow}
For every finite model satisfying Section~\ref{sec:model48},
\begin{equation}\label{eq:duality48}
 \sup_{\lambda\ge0}\sum_i\mu_i u_{0i}(\lambda)=\mathcal O(\mu).
\end{equation}
Both sides have finite attained optima. Adding variables $b_{ti}\in[0,\varepsilon]$, $1\le t<T$, and the equations
\begin{equation}\label{eq:shared48}
 \alpha_{tij}=f_{tij}b_{t+1,j},\qquad
 \gamma_{tij}=f_{tij}(\varepsilon-b_{t+1,j})
\end{equation}
to \eqref{eq:occupation48}--\eqref{eq:edge48} makes its optimum equal to the original all-restart common-policy value with objective measure $\mu$. Thus price optimization solves an explicitly identified relaxation; restoring shared successor budgets solves the stated policy problem.
\end{theorem}
\begin{proof}
Introduce free variables $v_{ti}$ and nonnegative variables $\lambda_{ti},h_{tij}$, with $v_T=\lambda_T=0$. The price linear program maximizes $\sum_i\mu_i v_{0i}$ subject to
\begin{align}
 v_{ti}-\beta\sum_jP_{tia,j}v_{t+1,j}
 +(\varepsilon-d_{tia})\lambda_{ti}
 -\beta\varepsilon\sum_jP_{tia,j}h_{tij}&\le k_{tia},\label{eq:priceLP48}\\
 h_{tij}\le\lambda_{ti},\qquad h_{tij}&\le\lambda_{t+1,j}.\nonumber
\end{align}
There are no $h$ variables in the last period. Given a price field, increasing $h$ to the corresponding minimum can only relax these inequalities. Backward maximization of $v$ then gives exactly \eqref{eq:price48}. Conversely, \eqref{eq:price48} and those minima are feasible. This identifies the linear-program value with the supremum on the left of \eqref{eq:duality48}.

Assign dual multipliers $y$ to \eqref{eq:priceLP48} and $\alpha,\gamma$ to the two upper bounds on $h$. The free $v$ variables give occupation conservation. The nonnegative $\lambda$ variables give \eqref{eq:budget48}; the nonnegative $h$ variables give \eqref{eq:edge48}. The resulting dual is precisely \eqref{eq:occupation48}. It is feasible: for any feasible policy set $y_{tia}=\beta^t\Pr_\mu^p(X_t=i)p_t(a\mid i)$, $\alpha_{tij}=f_{tij}D_{t+1,j}^p$, and $\gamma_{tij}=f_{tij}(\varepsilon-D_{t+1,j}^p)$. At a positive-date node the budget inequality is an equality; at a root it is its all-restart regret inequality. The primal price program is also feasible, for example with zero prices and sufficiently small $v$. Its value is bounded by a feasible policy's cost. Finite-dimensional linear-program duality proves equality and attainment.

The same policy construction satisfies \eqref{eq:shared48} with $b=D^p$, so the completed program has value at most the original optimum. In the reverse direction, take a completed feasible solution. At every node with $z_{ti}>0$, put $p_t(a\mid i)=y_{tia}/z_{ti}$. Conservation ensures that a positive-probability transition of this policy goes to a positive-occupation node. For $t\ge1$, substituting \eqref{eq:shared48} in \eqref{eq:budget48} and using incoming conservation gives
\[
 b_{ti}\ge\sum_a p_t(a\mid i)\{d_{tia}+\beta\sum_jP_{tia,j}b_{t+1,j}\},
\]
with terminal budget zero. At roots of positive $\mu$ mass the same argument replaces $b_{0i}$ by $\varepsilon$. Backward induction gives actual regret at most $b$ at positive-date occupied nodes and at most $\varepsilon$ at roots. At an unoccupied node choose a first maximizer of $r_t+\beta P_tJ_{t+1}^p$ after the later policy has been completed. If the future policy is globally feasible, this yields $J_t^p\ge V_t-\beta\varepsilon\ge V_t-\varepsilon$. Occupied transitions do not enter unoccupied successors, so this extension does not change their values. It is globally feasible and has objective $\sum ky$. This proves the reverse inequality.
\end{proof}

The theorem does not say that a usual occupation-measure formulation is incapable of expressing common policies. It gives the additional bilinear compatibility which this particular price relaxation omits. The problem-specific content is the precise edge allocation and its completion, not linear-program duality itself.

\begin{corollary}[Clipping and constant multipliers]\label{cor:clipping48}
In a finite model,
\begin{equation}\label{eq:masks48}
 \sup_\lambda L(\lambda)=\max_{S\subseteq\operatorname{supp}\nu}\mathcal O(\nu\ind_S).
\end{equation}
For a constant price $\lambda_{ti}=\ell$ at all nonterminal dates, the un-clipped objective is
\[
 \sum_i\nu_i u_{0i}=\inf_p\left\{\int C_0^p\,d\nu+\ell\left(\int D_0^p\,d\nu-\varepsilon\right)\right\}.
\]
Consequently the best constant-multiplier bound is no larger than the un-clipped state--date-price supremum, which is no larger than $\KR$.
\end{corollary}
\begin{proof}
For every vector $v$, $\sum_i\nu_i\max(0,v_i)=\max_S\sum_{i\in S}\nu_i v_i$. A finite maximum commutes with a supremum, so Theorem~\ref{thm:flow} proves \eqref{eq:masks48}. For the second statement, the price correction cancels successive constant-price allowance terms. Summing them along a trajectory leaves $-\ell\varepsilon$, while the disadvantages telescope to $D_0^p$. Backward minimization of the transformed objective proves the identity.
\end{proof}

The implementation uses a finite price cap and numerically proposed prices. It enumerates all nonempty masks when $n\le4$ and otherwise proposes the full mask and positive-mass singletons, subject to a declared time budget. Every proposed field is reconstructed rationally and ranked by its clipped certificate. A truncated mask search is not advertised as attaining \eqref{eq:masks48}; the earlier un-clipped-only objective is not silently equated with the displayed bound. Appendix~\ref{app:clipping48} gives an exact example in which the full-mask optimum is $3/8$ while a clipped field attains $15/32$.

\subsection{A strict separation with exact rational values}
\begin{theorem}\label{thm:separation}
There is a two-state, two-period, two-action model with action-independent transitions for which the constant-multiplier optimum, the clipped restart-price supremum, and the common-policy optimum are respectively
\begin{equation}\label{eq:strict48}
 \frac54\;<\;\frac{27}{16}\;<\;\frac74.
\end{equation}
The first two values have exact matching primal and dual witnesses. The final value follows from a one-variable affine minimization.
\end{theorem}
\begin{proof}
Take $\beta=\varepsilon=1/2$, zero terminal reward, and actions $G,B$ with operating rewards $0,-1$. Every transition goes to state 1. At date 0 the cost of $G$ is 1 in state 1 and 3 in state 2; at date 1 it is $3/2$ in either state. Action $B$ costs zero. Let $\nu=(1/2,1/2)$. Then $V=0$ and $d(G)=0,d(B)=1$.

Let $q$ denote the probability of $B$ at date 1 in state 1 and $q_i$ the two date-0 probabilities. Feasibility requires $0\le q\le1/2$ and $q_i\le1/2-q/2$. Costs decrease with each $q_i$, so they equal these upper limits at an optimum. The objective becomes $7/4+q/4$, minimized at $q=0$ with value $7/4$.

Set $\lambda_0=(1,3)$ and $\lambda_1=(3/2,3/2)$. The recurrence gives $u_1=(3/4,3/4)$ and $u_0=(9/8,9/4)$, hence the lower value $27/16$. A feasible occupation witness in \eqref{eq:occupation48} takes the date-0 probabilities of $B$ as $(1/4,1/2)$ and the date-1 probability as $1/4$. The discounted occupations are $y_{01}=(3/8,1/8)$, $y_{02}=(1/4,1/4)$, and $y_{11}=(3/8,1/8)$. On the two incoming edges use $(\alpha,\gamma)=(1/8,0)$ and $(0,1/8)$. Every occupation, node-budget, and edge inequality holds, and the objective is $27/16$. These matching witnesses prove the un-clipped supremum.

Clipping cannot improve it. With initial mass only at state 1 the original optimum is $9/8$ (choose $q=1/2$); with mass only at state 2 it is $9/4$ (choose $q=0$). Every price potential at the corresponding root is bounded by that positive optimum, hence so is its positive part. The same price field attains both bounds simultaneously.

Finally, for the expected-initial-regret problem choose $(q_1,q_2,q)=(0,1,0)$. Its average regret is $1/2$ and its cost is $5/4$. The constant price $\ell=3$ gives $u_0=(1/4,9/4)$ and the same value. Weak duality proves the constant-multiplier optimum $5/4$. The conflicting preferred terminal randomizations explain the gap between the middle and final values: incoming edges cannot deploy different values of $q$ at their common successor.
\end{proof}

The separation supplies a constructive reason to refine shared policy variables rather than assuming that more price proposals alone must close every interval. It also distinguishes a policy-to-certificate identity from a true optimization error. The former's clipping correction is explicitly negative: if its un-clipped nonnegative slack sum is $\mathcal S$, then $C^p(\nu)-L(\lambda)=\mathcal S-\int(-u_0)_+d\nu$. We do not call this clipped expression a sum of nonnegative terms. The full earlier identity and its case-level diagnostics are retained in the supplement.
